% !TITLE 春夜喜雨
% !AUTHOR 杜甫
% !DYNASTY 唐
% !ORDER 25

\begin{poem}
好雨知时节，当春乃发生。
随风潜入夜，润物细无声。
野径云俱黑，江船火独明。
晓看红湿处，花重锦官城。
\end{poem}

\begin{appreciation}
这首诗写于上元二年（761年）的春天，杜甫刚到成都不久。那一年成都大旱，百姓焦急地等待着春雨，终于在一个夜晚，雨来了。

好雨知时节，当春乃发生——好雨知道时节，正当春天万物需要生长的时候就来了。知字是全诗的灵魂。雨怎么会知道时节呢？但在这场春雨面前，你会觉得它真的知道。它知道庄稼在等着它，知道花在等着它，知道人在等着它。这种拟人不是修辞技巧，是人在极度渴望之后的感恩心理。

随风潜入夜，润物细无声——随着风在夜里悄悄潜入，滋润万物却一点儿声音也没有。潜字写得好：雨不是大摇大摆来的，是悄悄来的，是体贴的，是不打扰人的。细无声更妙：它不是没有声音，而是声音太细了，细到听不见。这种细微，正是春雨的特征，也是杜甫观察力的体现。

野径云俱黑，江船火独明——田野间的小路一片漆黑，只有江上的渔船还亮着一点灯火。这是一幅深夜雨中的画面：天地之间几乎全黑，只有那一点渔火在闪烁。俱黑和独明形成了强烈的对比，那一点灯火在黑暗中显得尤为珍贵，像希望一样。

晓看红湿处，花重锦官城——天亮时看那被雨水打湿的红花，锦官城里繁花沉甸甸的。重字是这首诗最精妙的一个字。花因为沾了雨水而变重，低垂下来。这是一个视觉的细节，也是一个触觉的想象：你能感觉到那些花瓣上水珠的分量。锦官城就是成都。杜甫刚到成都，对这个城市还充满新鲜感和感激——感激它给了自己一个避难之所，感激这场春雨给了这座城市生机。
\end{appreciation}
